\section{Research design and relation to the papers}

The experiment reproduces the mixed-copula procedure of Sabino da Silva,
Ziegelmann and Caldeira (Paper~1), with the marked-to-market capital conventions
of Rad, Low and Faff (2016, Paper~2). All page references below denote physical
pages of the supplied PDFs, including their cover pages. Paper~1 studies historical
S\&P~500 constituents, allows mixtures to represent different dependence patterns,
and uses six-month trading windows. Paper~2 compares distance, cointegration and
single-copula strategies across the CRSP ordinary-share universe, emphasizing
trading costs and capital denominators. Their different universes, models and
portfolio construction make their empirical rankings distinct research findings,
rather than contradictory predictions for these three pairs.

\begin{table}[htbp]
\centering
\small
\textbf{Reproduction map and declared adaptations.}\par\medskip
\begin{tabularx}{\linewidth}{@{}>{\raggedright\arraybackslash}p{0.15\linewidth}>{\raggedright\arraybackslash}X>{\raggedright\arraybackslash}X>{\raggedright\arraybackslash}X@{}}
\toprule
Component & Paper 1 & This experiment & Paper 2 \\
\midrule
Universe & Historical S\&P~500 members; 1990--2015 data & AAPL--GOOG, IBM--SPY, DIA--SPY; 2006--2025 evaluation & CRSP ordinary shares; 1962--2014 \\
Windows & 12 months formation; 6 months trading; six-month advance & Same calendar lengths; January/July starts & Monthly starts with six overlapping portfolios \\
Pair selection & Rank normalized-price SSD; portfolios of 5--35 pairs & Three prescribed pairs retained in fixed capital slots & SSD selection; 20 pairs; extra cointegration screen \\
Dependence & Five single copulas and CFG/CtG mixtures; likelihood selection & Same seven candidates; BIC and best-single controls & Five single families, including survival Clayton/Gumbel; information criteria \\
Marginals & ARMA--GARCH standardized residuals and empirical CDF & Four ARMA orders; Gaussian QMLE and AIC made explicit & Parametric distributions fitted to daily returns \\
Execution and PnL & Same-day closing prices; committed/fully invested measures & Next-close execution; fixed units; monthly RCC/REC & Marked-to-market PnL; \$1 long-side denominator \\
\bottomrule
\end{tabularx}
\end{table}

The data cutoff is December 2025, the final complete calendar year selected before
model fitting. The first evaluation year is 2006, with the preceding year available
for formation. Parameters are estimated using the preceding 12 calendar months
and held fixed during each following six-month January--June or July--December
window. This follows Paper~1's calendar construction (p.~6); Paper~2's six
overlapping monthly cohorts are not implemented (p.~5). Each pair has an independent
capital allocation, including the two uses of SPY. Positions and transaction costs
are not netted across pairs. Cointegration rejections leave their allocation in cash.

This is a methodological reproduction using the assignment's suggested pairs.
It does not reconstruct historical index membership or the original SSD-ranked
universe. The ETFs also fall outside Paper~2's ordinary-share restriction. Present-day
pair selection, the later sample, and three dependent pair series limit conclusions
about either original universe (Paper~1, p.~11; Paper~2, p.~5).

\section{Marginal models, copulas and trading signals}

\subsection{Formation-only estimation}

For adjusted total-return marks $P_{j,t}$, define log returns
$r_{j,t}=\log(P_{j,t}/P_{j,t-1})$. Sklar's representation separates marginal
distributions from dependence:
\begin{equation}
 H(x,y)=C\{F_1(x),F_2(y)\},\qquad
 f(x,y)=c\{F_1(x),F_2(y)\}f_1(x)f_2(y).
\end{equation}
Thus nonnormal marginal behavior can be modeled separately from joint tail
behavior (Paper~1, pp.~7--8; Paper~2, p.~7).

Each return series has an ARMA$(p,q)$--GARCH$(1,1)$ marginal, with
$p,q\in\{0,1\}$:
\begin{align}
 r_t&=\mu+\phi(r_{t-1}-\mu)+\theta\epsilon_{t-1}+\epsilon_t,
 &\epsilon_t&=\sigma_t z_t,\\
 \sigma_t^2&=\omega+\alpha\epsilon_{t-1}^2+\beta\sigma_{t-1}^2.
\end{align}
Absent AR or MA terms have coefficient zero. Gaussian quasi-maximum likelihood
estimates each candidate, and minimum AIC selects among the four orders.
Paper~1 specifies the order range and residual transformation (pp.~9--10), but
leaves this likelihood and order-selection convention open; they are explicit
implementation choices. Positive variance and stationary variance-recursion
constraints are imposed.

The formation standardized residuals $\widehat z_t$ give pseudo-observations
\begin{equation}
 u_t=\frac{n}{n+1}\widehat F_n(\widehat z_t)
     =\frac{\#\{s:\widehat z_s\leq\widehat z_t\}}{n+1}.
\end{equation}
During trading, coefficients and the empirical CDF remain frozen; observed returns
update only the causal mean/variance recursion. Future residuals are mapped through
that formation CDF and clipped to $[1/(n+1),n/(n+1)]$. Neither the CDF nor parameter
selection uses future observations.

The candidate set is Clayton, Frank, Gumbel, Gaussian, Student-$t$, and mixtures
CFG and CtG. For components $k\in\{1,2,3\}$,
\begin{equation}
 C_{\mathrm{mix}}(u,v)=\sum_k\pi_k C_k(u,v;\vartheta_k),\qquad
 (\widehat\pi,\widehat\vartheta)=\arg\max_{\pi,\vartheta}
 \sum_{t\in\mathrm{formation}}\log\left[\sum_k\pi_k
 c_k(u_t,v_t;\vartheta_k)\right],
\end{equation}
where $\pi_k\geq0$ and $\sum_k\pi_k=1$. CFG combines Clayton, Frank and Gumbel;
CtG replaces Frank with Student-$t$ (Paper~1, p.~10). The logarithm applies to
the weighted density sum, rather than to separate weighted log densities.
The primary strategy chooses maximum unpenalized formation log likelihood, as
in Paper~1 (p.~12). Because flexible mixtures have an in-sample advantage, minimum
$\mathrm{BIC}=k\log n-2\ell$ and maximum likelihood restricted to the five single
copulas provide controls. These are predetermined selection rules, not rankings
chosen from trading-period profits.

\subsection{Conditional ranks and position rules}

Conditional distributions and accumulated flags are
\begin{align}
 h_{1,t}&=\frac{\partial C(u_t,v_t)}{\partial v},&
 h_{2,t}&=\frac{\partial C(u_t,v_t)}{\partial u},\\
 M_{j,t}&=M_{j,t-1}+h_{j,t}-\tfrac12,&M_{j,0}&=0.
\end{align}
Mixture conditionals are the same weighted sums of component conditionals.
An extreme conditional rank identifies an unusual contemporaneous residual
relative to its partner. It is not a calibrated probability of a future price
reversal, and accumulating ranks does not establish stationarity.

While flat, $M_1>0.2$ and $M_2<-0.2$ trigger short first/long second; the
opposite inequalities trigger long first/short second. The primary exit occurs
when \emph{either} flag reaches or crosses zero from its entry-side direction.
This follows Paper~1's numbered step~5 (p.~11), although its preceding prose
says both flags (p.~9). A separate sensitivity requires both flags to be on
their converged sides simultaneously, consistent with Paper~2's wording
(p.~8). Flags reset only at each six-month window's start. No same-signal-day
exit/re-entry is allowed.

Signals observed at close $t$ execute at close $t+1$; a new position earns no
return before that execution. This is a declared departure from Paper~1's
same-close assumption (p.~11). Outstanding positions are liquidated at the
predefined final close, and terminal-date entries are suppressed.

\subsection{Distance and cointegration controls}

The distance control scales formation total-return indices to one and estimates
their spread standard deviation. It rebases prices at trading start, enters at
two formation standard deviations and exits at a zero crossing, following
Paper~2 (p.~6). The cointegration control fits normalized second price on first,
with an intercept, and applies an Engle--Granger test at 5\%, using AIC lag
selection. Eligible pairs require positive $\beta$; the centered formation spread
$P_2-\beta P_1$ is standardized and traded at $\pm2$ with zero-crossing exits.
Below $-2$, the allocation is long $K$ of asset~2 and short $\beta K$ of asset~1;
above $+2$, it is long $K$ of asset~1 and short $K/\beta$ of asset~2. This follows
Paper~2's operational \emph{dollar} rule (p.~7), rather than treating $\beta$ as
an invariant share ratio. Rejected pairs are not replaced from a larger universe.

\section{PnL accounting, costs and capital}

Let $K=\$100{,}000$ be the fixed long-side stake per pair. For long asset~1 and
short asset~2 entered at close $e$,
\begin{equation}
 q_1=\frac{K}{P_{1,e}},\qquad q_2=-\frac{K}{P_{2,e}},\qquad
 \mathrm{PnL}^{\mathrm{gross}}_t=
 \sum_{j=1}^2q_{j,t-1}(P_{j,t}-P_{j,t-1}).
\end{equation}
Units remain fixed until exit. Equivalently, writing simple returns as $R_{j,t}$,
\begin{equation}
 \frac{\mathrm{PnL}^{\mathrm{gross}}_t}{K}
 =\sum_j\underbrace{\frac{q_{j,t-1}P_{j,t-1}}{K}}_{w_{j,t-1}}
 R_{j,t}.
\end{equation}
The weights drift with marked prices; holding them constantly at $+1,-1$
would introduce daily rebalancing. Adjusted-price units are synthetic total-return
units, so dividends are not separately credited. The equal-dollar pair starts
with $2K$ gross exposure but has denominator $K$, matching Paper~2's convention
(pp.~8--9); gross-exposure returns would be half as large initially. The three-slot
portfolio commits $3K$. Idle cash earns zero, and profits are not reinvested into
larger subsequent stakes.

At each fill, proportional transaction cost and optional daily borrow cost are
\begin{equation}
 \mathrm{TC}_t=c\sum_j|\Delta q_{j,t}|P_{j,t},\qquad
 \mathrm{BC}_t=\frac{b}{252}\sum_j\max(-q_{j,t-1},0)P_{j,t-1}.
\end{equation}
Net PnL subtracts both, including liquidation costs. The base $c=5$ basis
points per leg per fill implies 20 basis points per complete equal-dollar pair
at unchanged prices, matching Paper~1's round-trip benchmark (pp.~11--12).
The fee grid is 0, 1, 5, 10, 14.5 and 20 basis points; borrow scenarios are
0\%, 1\% and 3\% annually. These are scenarios, not reconstructed historical
quotes. Paper~2 uses a historical commission/impact schedule (p.~8); the
14.5-basis-point case corresponds to its stated recent 58-basis-point pair
round trip only when notionals remain unchanged.

For monthly pair return $r_{i,m}=\sum_{t\in m}\mathrm{PnL}^{\mathrm{net}}_{i,t}/K$,
\begin{equation}
 \mathrm{RCC}_m=\frac{\sum_{i=1}^3 r_{i,m}}{3},\qquad
 \mathrm{REC}_m=\frac{\sum_{i=1}^3r_{i,m}}{n_m}.
\end{equation}
Here $n_m$ counts pairs held or traded during the month, including carry-in
positions. REC is undefined when none trade; RCC is then zero. These adapt
Paper~2's equations~(14)--(15), p.~8, to three nonoverlapping slots. PnL is
recognized daily, including unrealized gains, rather than assigned wholly to
exit months. REC is an activity-conditioned reporting measure; committed capital
underlies the portfolio equity curve. Equity equals initial capital plus cumulative
PnL, with no geometric compounding of fixed-stake returns.

Accounting checks use hand-computed paths to verify execution lag, drifting
weights, unchanged-price round-trip fees, forced liquidation, borrow charges,
hedge-ratio timing, and equality of daily and completed-trade net PnL totals.
